% Part: applied-modal-logics
% Chapter: temporal-logic
% Section: possible-histories
%READERNOTE{SOL6-A187: इतिहासात कालिक स्थान स्पष्ट होण्यासाठी प्रत्येक अवस्था-घटना स्वतंत्र मानली; प्रणालीची स्थिती पुन्हा आली तरी तिच्या वेगळ्या वेळच्या घटना वेगळ्या आहेत. Boolean उपसूत्रे त्याच बिंदू व इतिहासावर तपासायची हे स्पष्ट केले. शेवटच्या उदाहरणातील कधीच पुढे सत्य होत नाही हा सार्विक नकार दुरुस्त केला.}

\documentclass[../../../include/open-logic-section]{subfiles}

\begin{document}

\olfileid{aml}{tl}{poss}

\olsection{संभाव्य इतिहास}

आतापर्यंत वापरलेली कालिक तर्कशास्त्राची संबंधी
प्रतिमाने अत्यंत लवचिक आहेत, कारण प्राप्यता संबंधावर
कोणतेही निर्बंध घालावे लागत नाहीत. त्यामुळे कालिक
प्रतिमाने भूतकाळात आणि भविष्यकाळात दोन्हीकडे
शाखायुक्त असू शकतात. पण शाखा फुटण्याची अधिक
``मोडल'' कल्पना आपण विचारात घेऊ शकतो: घटनांचे
अनुक्रम हे संभाव्य इतिहास मानायचे. यासाठी भाषा
बदलणे आवश्यकच नाही; तरीही आपण ``नेहमीचे''
मोडल संकारक $\Box$ आणि~$\Diamond$ जोडू शकतो.
कालांतराने कर्त्यांच्या ज्ञानात होणारे बदल दर्शवण्यासाठी
ज्ञानविषयक प्राप्यता संबंध जोडण्याचाही विचार करता येतो.

\begin{defn}
  कालिक भाषेसाठीचे \emph{संभाव्य इतिहासांचे प्रतिमान}
  हे $\mModel{M} = \tuple{T, C, V}$ असे त्रिक असते. येथे
  \begin{enumerate}
    \item $T$ हा रिक्तेतर संच असून त्यातील घटकांचा अर्थ
      काळातील अवस्था असा लावला जातो.
    \item $C$ हा प्रणालीच्या संगणकीय मार्गांचा, म्हणजे
      \emph{संभाव्य इतिहासांचा}, संच आहे. दुसऱ्या शब्दांत,
      $C$ हा $s_1$, $s_2$,~$s_3$, \dots अशा अवस्थांच्या
      $\sigma$ या अनुक्रमांचा संच आहे; येथे
      प्रत्येक $s_i \in T$. येथे $s_i$ या काळातील स्वतंत्र
      अवस्था-घटना मानतो: एका इतिहासात एकाच $T$-घटकाची
      पुनरावृत्ती होत नाही. प्रणालीची तीच स्थिती पुन्हा आली,
      तरी तिच्या वेगवेगळ्या काळातील घटनांना वेगळे घटक मानायचे.
      त्यामुळे $t,\sigma$ याने त्या इतिहासातील एकच स्थान
      निश्चित होते आणि $\prec_\sigma$ चा पुढील अर्थ संदिग्ध राहत नाही.
    \item $V$ हे प्रत्येक !!{propositional
      variable}~$p$ ला कालबिंदूंचा संच~$V(p) \subseteq T$ देणारे फलन आहे.
  \end{enumerate}
  गोष्टी सोप्या करण्यासाठी, एखादा इतिहास~$C$ मध्ये
  असेल, तर त्याचे सुरुवातीचे घटक वगळून मिळणारे
  सर्व उर्वरित अनुक्रमही त्यात आहेत, असे आपण साधारणतः
  मानू. उदाहरणार्थ, $s_1$, $s_2$,~$s_3$ हा अनुक्रम~$C$
  मध्ये असेल, तर $s_2$, $s_3$ आणि~$s_3$ हेही त्यात
  असतात. तसेच, अवस्था $s_i$ आणि~$s_j$ या
  अनुक्रम~$\sigma$ मध्ये असतील, तर $i < j$ असताना
  $s_i \prec_\sigma s_j$ असे म्हणतो.
  $t \in V(p)$ असेल, तर $p$ हे~$t$ येथे
  \emph{सत्य आहे} असे म्हणतो.
\end{defn}

येथे संबंधित बदल असा आहे: प्रतिमान~$\mModel M$
मध्ये कालबिंदू~$t$ येथे !!a{formula} ची सत्यता
तपासताना आपण ती $t$ ही अवस्था-घटना असलेल्या
इतिहास~$\sigma \in C$ च्या सापेक्ष तपासतो.
!!{propositional variable}s ची सत्यता फक्त $t$ वर अवलंबून असते:
$\mSat{M}{p}[t,\sigma]$ जर आणि फक्त जर $t\in V(p)$.
सत्यता-फलनात्मक संयोजकांच्या अटी
\olref[sem]{defn:tmodels} मधीलप्रमाणेच राहतात;
मात्र त्यांतील प्रत्येक उपसूत्राची सत्यता त्याच
$t,\sigma$ येथे तपासायची. उदाहरणार्थ,
$\mSat{M}{\lnot !B}[t,\sigma]$ जर आणि फक्त जर
$\mSat/{M}{!B}[t,\sigma]$; आणि
$\mSat{M}{!B\land !C}[t,\sigma]$ जर आणि फक्त जर
$\mSat{M}{!B}[t,\sigma]$ आणि $\mSat{M}{!C}[t,\sigma]$.
संयोजक स्वतः इतिहास बदलत नसले, तरी त्यांची कालिक किंवा
मोडल उपसूत्रे त्या इतिहासावर अवलंबून असू शकतात.
आता या इतिहासांच्या सापेक्ष आपण
भविष्यकाळी संकारक~$\Ftemp$ ची पुनर्व्याख्या करतो
आणि~$\Diamond$ संकारक जोडतो.

\begin{defn}\ollabel{defn:phmodels}
  प्रतिमान~$\mModel M$ मध्ये~$t, \sigma$ येथे
  !!a{formula}~$!A$ ची \emph{सत्यता}, म्हणजेच
  $\mSat{M}{!A}[t, \sigma]$:
  \begin{enumerate}
  \item\ollabel{defn:sub:phmodels-f} \indcase{!A}{\Ftemp
    !B}{$\mSat{M}{\indfrm}[t, \sigma]$ जर आणि फक्त जर
    असा एखादा $t' \in T$ असेल की $t \prec_\sigma t'$
    आणि $\mSat{M}{!B}[t', \sigma]$.}
  \item\ollabel{defn:sub:phmodels-diamond} \indcase{!A}{\Diamond
    !B}{$\mSat{M}{\indfrm}[t, \sigma]$ जर आणि फक्त जर
    असा एखादा $\sigma' \in C$ असेल की त्यात $t$ येतो
    आणि $\mSat{M}{!B}[t, \sigma']$.}
  \end{enumerate}
\end{defn}

इतर कालिक आणि मोडल संकारकांचीही अशीच व्याख्या
करता येते. आता आपण काळवाचकता आणि मोडलता
एकत्र करणारी विधाने दर्शवू शकतो. उदाहरणार्थ,
``$p$ पुढे घडणार नाही, पण घडू शकला असता'' हे
विधान $\lnot \Ftemp p \land \Diamond \Ftemp p$
या सूत्राने दर्शवता येईल. ज्या बिंदू $t$ आणि इतिहास $\sigma$ येथे
त्या इतिहासातील कोणत्याही पुढील अवस्थेत $p$ सत्य नसते,
पण $t$ असलेल्या एखाद्या पर्यायी इतिहासात $p$ पुढील
एखाद्या अवस्थेत सत्य होते, तेथे हे सूत्र सत्य ठरेल.

\end{document}
